% Generated draft: requires mathematical and source review.
\documentclass[11pt]{article}
\usepackage[margin=25mm]{geometry}
\usepackage{fontspec}
\setmainfont{Latin Modern Roman}
\usepackage{amsmath,amssymb,mathtools,unicode-math}
\setmathfont{Latin Modern Math}
\usepackage{enumitem,needspace,xurl,hyperref}
\setlength{\parindent}{0pt}
\setlength{\parskip}{5pt}
\setlength{\emergencystretch}{3em}
\newcommand{\N}{\mathbb{N}}
\newcommand{\Z}{\mathbb{Z}}
\newcommand{\Q}{\mathbb{Q}}
\newcommand{\R}{\mathbb{R}}
\newcommand{\C}{\mathbb{C}}
\newcommand{\F}{\mathbb{F}}
\newcommand{\A}{\mathbb{A}}
\newcommand{\PP}{\mathbb{P}}
\newcommand{\E}{\mathbb{E}}
\newcommand{\eps}{\varepsilon}
\newcommand{\dd}{\,\mathrm{d}}
\begin{document}
\section*{UnsolvedMath ID: 800060}
\textbf{Source problem code:} catalogue entry
\subsection{Definitions and mathematical statement}
The supplied record calls this the \emph{Virtual Euler class one conjecture}. Its main formula suggests the following formulation, conditional on orientation, parity, and coefficient conventions that the source does not fully specify.

Assume that \(M\) is a closed oriented hyperbolic \(3\)-manifold, and write \(b_1(M)=\dim_{\R}H^1(M;\R)\). For an integral homology class \(z\in H_2(M;\Z)\), the Thurston norm is
\[
 x_M(z)=\min_{[S]=z}\sum_i\max\{0,-\chi(S_i)\},
\]
where the minimum ranges over embedded closed oriented surfaces \(S=\bigsqcup_i S_i\) representing \(z\). Extend this seminorm to \(H_2(M;\R)\). Its dual on \(H^2(M;\R)\) is
\[
 x_M^*(a)=\sup\bigl\{|\langle a,z\rangle|:z\in H_2(M;\R),\ x_M(z)\leq 1\bigr\}.
\]
Let \(\iota_X:H^2(X;\Z)\to H^2(X;\R)\) denote the coefficient homomorphism. An integral real class belongs to \(\operatorname{im}\iota_X\); the conventional reading of an even integral real class is membership in \(\operatorname{im}(2\iota_X)\).

A codimension-one foliation is taut if every leaf meets a closed curve transverse to the foliation. For a foliation \(\mathcal F\) whose tangent plane bundle is oriented, let
\[
 e_{\Z}(T\mathcal F)\in H^2(X;\Z),\qquad
 e_{\R}(\mathcal F)=\iota_X\bigl(e_{\Z}(T\mathcal F)\bigr)
\]
be its integral and real Euler classes.

\emph{Question, under these conventions.} For every such \(M\) with \(b_1(M)\geq 1\), and every
\[
 a\in\operatorname{im}(2\iota_M),\qquad x_M^*(a)=1,
\]
does there exist a finite covering map \(p:\widetilde M\to M\) and a taut codimension-one foliation \(\mathcal F\) of \(\widetilde M\), with oriented tangent plane bundle, such that
\[
 e_{\R}(\mathcal F)=p^*a
 \quad\text{in }H^2(\widetilde M;\R)?
\]
Here the displayed formulation uses the ordinary pullback \(p^*:H^2(M;\R)\to H^2(\widetilde M;\R)\).

\emph{Unresolved source qualifications.} The source says ``(even) integral'' without explaining the parentheses, and does not explicitly require orientability of \(M\), an orientation of the foliation's tangent bundle, or coorientability of the foliation. The conventions above are therefore conditional interpretations requiring review. Moreover, the source instead describes \(p^*\) as acting on \(H^1\) and claims Poincar\'e duality identifies \(H^1\) with \(H^2\). For a closed oriented \(3\)-manifold, Poincar\'e duality identifies \(H^1\) with \(H_2\) and \(H^2\) with \(H_1\), not canonically \(H^1\) with \(H^2\). Thus that explanation does not specify the asserted identification or its compatibility with pullback. Finally, since \(a\) is supplied only as a real cohomology class, the displayed equality is interpreted in real cohomology; an intended integral equality would require specification of an integral lift and its torsion information.
\subsection{Short English statement}
Under the stated conditional conventions, the question asks whether every even integral real cohomology class of dual Thurston norm exactly one on a closed oriented hyperbolic \(3\)-manifold with positive first Betti number becomes the Euler class of a taut foliation after passage to a finite cover. The source's orientation, parity, coefficient, and pullback conventions require clarification.
\subsection{Sources}
\begin{itemize}
\item Supplied local UnsolvedMath registry record: permanent ID \(800060\), exact source code \texttt{catalogue entry}; title: \emph{Foliations and Diffeomorphism Groups --- Question 1 (Virtual Euler class one conjecture)}. The record attributes the problem to Nariman and Yazdi, with proposed year \(2025\).
\item Nariman and Yazdi, \emph{Problem List on Foliations and Diffeomorphism Groups} (2025), Question 1. Supplied source URL: \url{https://arxiv.org/abs/2503.18274}. The record reports access on 2026-10-16 and extraction by \texttt{arxiv\_latex}; these are supplied provenance details, not a claim of independent source access or verification.
\item The local record labels the problem \texttt{open} and gives an update timestamp of \texttt{2026-10-16T14:00:00Z}. This status is supplied dated curation, not an independently verified literature assessment. No additional literature references are supplied.
\end{itemize}
\end{document}
